Os sistemas autónomos e os sistemas embebidos integram funções de
perceção, processamento, controlo e atuação que dependem da troca de
informação entre os seus elementos. A aquisição de dados de sensores,
o envio de comandos a atuadores e a comunicação com controladores ou
equipamentos de supervisão podem envolver tecnologias de comunicação
distintas numa mesma instalação. Esta coexistência resulta das
características dos dispositivos e das funções que desempenham, bem
como da necessidade de integrar equipamentos desenvolvidos em momentos
ou contextos diferentes.

Neste contexto, sensores, atuadores, controladores e outros equipamentos
podem disponibilizar interfaces como I2C, RS232, RS485, CAN e Ethernet.
Estas tecnologias apresentam diferenças nas características elétricas,
na organização das ligações e nos mecanismos de comunicação que
suportam. Por conseguinte, a ligação de um novo equipamento exige a
análise da compatibilidade da interface e da forma como os dados são
transmitidos e recebidos, além da compreensão do seu significado para
a aplicação.

O problema central reside na heterogeneidade destas interfaces e na
dificuldade de integrar operações de leitura e escrita sobre equipamentos
com formas de acesso distintas. O suporte de cada tecnologia pode exigir
configurações, procedimentos de comunicação e mecanismos de tratamento
de erros específicos. Quando estas particularidades são tratadas de
forma dispersa na aplicação, a inclusão ou substituição de equipamentos
pode implicar alterações em várias partes do sistema, dificultando a
manutenção e a reutilização das funções de comunicação. Além disso, a
disponibilidade de uma ligação compatível não assegura, por si só, a
interpretação correta dos dados trocados.

A relevância do trabalho reside, assim, na possibilidade de organizar de
forma coerente os mecanismos de acesso aos equipamentos, facilitando a
integração e a evolução de sistemas com diferentes tecnologias de
comunicação.
